% JETSPIN compiling documentation

\section{Compiling the Source Code\label{sec:Compiling-the-Source}}

The \textit{build} sub-directory stores a UNIX \textbf{Makefile} that
assembles serial and parallel executable versions with different
compilers. JETSPIN may be compiled on any UNIX platform. From the root
of the repository, invoke this file directly while making \textit{source}
the working directory; there is no need to copy the Makefile. Use

\bigskip{}

\texttt{make -C source -f ../build/Makefile target}

\bigskip{}

where \texttt{target} is one of the options reported in
Table~\ref{Tab:targets}. Object and module files are generated in
\textit{source}, and the resulting \texttt{main.x} executable is placed
in \textit{execute}. For example, the serial GFortran build is produced
with

\bigskip{}

\texttt{make -C source -f ../build/Makefile gfortran}

\bigskip{}

and the corresponding Open MPI build with target
\texttt{gfortran-mpi}. Build products can be removed with target
\texttt{clean} using the same command form.

NVIDIA HPC SDK provides separate CPU-only and OpenACC build targets. On a
module-based installation, first load the SDK and verify that its MPI wrapper
uses NVFORTRAN:

\begin{verbatim}
module purge
module use /opt/nvidia/hpc_sdk/modulefiles
module load nvhpc/24.3
nvfortran --version
mpif90 --showme:command
\end{verbatim}

The serial and MPI CPU-only builds are selected respectively with
\texttt{nvfortran} and \texttt{nvfortran-mpi}. The latter requires
\texttt{mpif90 --showme:command} to report \texttt{nvfortran}.
Neither target enables OpenACC offloading.

The \texttt{nvfortran-openacc} target enables single-GPU OpenACC
offloading. Its default values target compute capability 8.0, including
NVIDIA A30 and A100 GPUs, and the CUDA 12.3 toolkit distributed with HPC SDK
24.3. It also selects the \texttt{nofma} GPU option to preserve the numerical
behaviour of the nearly straight three-point curvature calculation:

\begin{verbatim}
make -C source -f ../build/Makefile nvfortran-openacc
\end{verbatim}

The OpenACC GPU and host-validation targets expose the standard
\texttt{\_OPENACC} preprocessing macro.
Accelerator imports, initialization, Coulomb dispatch, and equation-of-motion
dispatch are enclosed by this macro. CPU targets preprocess the same sources
without defining it, so accelerator calls are absent from their compiled
translation units rather than being retained as no-op runtime branches.

Another device architecture or installed CUDA toolkit can be selected through
the \texttt{GPUCC} and \texttt{CUDA\_VERSION} Make variables. The
compute capability is passed without the \texttt{cc} prefix; for example:

\begin{verbatim}
make -C source -f ../build/Makefile nvfortran-openacc \
  GPUCC=90 CUDA_VERSION=12.3
\end{verbatim}

Target \texttt{nvfortran-openacc-host} compiles the OpenACC control path
for host execution. It is intended for compiler and numerical validation when
an NVIDIA device is unavailable and is not a GPU performance build. Before a
device run, \texttt{nvidia-smi} and \texttt{nvaccelinfo} should both
confirm that the selected GPU is visible.

Target \texttt{nvfortran-openacc-host-forces} is a development-only GPU build
that evaluates the evaporative Coulomb and Maxwell geometric forces with the
trusted host routines and copies their results back to the device. It is useful
for numerical isolation, but the per-stage transfers make it unsuitable for
performance measurements or production simulations.

On Windows system we advice the user to compile JETSPIN under the
command-line interface Cygwin \cite{racine2000cygwin}. Finally, the
binary executable file can be run in the \textit{execute} sub-directory.
Note that a FORTRAN 90 compiler is necessary in order to compile JETSPIN.
Further, a Message Passing Interface Implementation should be also
present, if the user wishes to compile JETSPIN in parallel mode. We
recommend GFORTRAN as FORTRAN 90 compiler, but IFORT-Intel(R) Fortran
Compiler or G95 can also be used. As Message Passing Interface Implementation,
we suggest to use that provided by the Open MPI Project.

\begin{table}[H]
\begin{centering}
\begin{tabular}{p{0.29\textwidth}p{0.65\textwidth}}
\hline
\textbf{target:}  & \textbf{meaning:}\tabularnewline
\hline
\hline
gfortran  & compile in serial mode using the GFortran compiler.\tabularnewline
gfortran-mpi  & compile in parallel mode using the GFortran compiler and the Open
Mpi library.\tabularnewline
cygwin  & compile in serial mode using the GFortran compiler under the command-line\tabularnewline
 & interface Cygwin for Windows.\tabularnewline
cygwin-mpi  & compile in parallel mode using the GFortran compiler and the Open
Mpi library\tabularnewline
 & under the command-line interface Cygwin for Windows (note a precompiled\tabularnewline
 & package of the Open Mpi library is already available on Cygwin).\tabularnewline
intel  & compile in serial mode using the Intel compiler.\tabularnewline
intel-mpi  & compile in parallel mode using the Intel compiler and the Intel Mpi
library.\tabularnewline
intel-openmpi  & compile in parallel mode using the Intel compiler and the Open Mpi
library.\tabularnewline
nvfortran & compile in serial CPU mode using NVIDIA Fortran; OpenACC offloading
is disabled.\tabularnewline
nvfortran-mpi & compile in parallel CPU mode using an MPI wrapper configured
with NVIDIA Fortran; OpenACC offloading is disabled.\tabularnewline
nvfortran-openacc & compile a single-GPU OpenACC executable; configure the
device with \texttt{GPUCC} and \texttt{CUDA\_VERSION}.\tabularnewline
nvfortran-openacc-host & compile the OpenACC code path for host-only validation.\tabularnewline
nvfortran-openacc-host-forces & compile the GPU path with development-only
host force evaluation and per-stage transfers.\tabularnewline
help  & return the list of possible target choices\tabularnewline
\hline
\end{tabular}
\par\end{centering}
\protect\caption{List of targets for several common workstations and parallel computers,
which can be used with the Makefile in the \textit{build} sub-directory.}

\label{Tab:targets}
\end{table}
